\documentclass[11pt]{article}

\usepackage[margin=1.1in]{geometry}
\usepackage{amsmath,amssymb,amsthm,mathtools}
\usepackage{enumitem}
\usepackage{hyperref}
\usepackage{array}
\usepackage[title]{appendix}
\usepackage[nameinlink,noabbrev,capitalise]{cleveref}

\hypersetup{colorlinks=true,linkcolor=blue,citecolor=blue,urlcolor=blue}

\newtheorem{theorem}{Theorem}[section]
\newtheorem{lemma}[theorem]{Lemma}
\newtheorem{corollary}[theorem]{Corollary}
\theoremstyle{remark}
\newtheorem{remark}[theorem]{Remark}

\providecommand{\doi}[1]{\href{https://doi.org/#1}{doi:#1}}

\newcommand{\E}{\mathbb{E}}
\newcommand{\Prob}{\mathbb{P}}
\newcommand{\R}{\mathbb{R}}
\newcommand{\C}{\mathbb{C}}
\newcommand{\Var}{\operatorname{Var}}

\title{The right tail of higher-order stochastic Airy operators and edges of unitary ensembles}
\author{Yitzchak Shmalo\\
Einstein Institute of Mathematics, The Hebrew University of Jerusalem\\
Givat Ram, Jerusalem, Israel}
\date{}

\begin{document}
\maketitle

\begin{abstract}
Krishnapur, Rider and Vir\'ag proposed a random Schr\"odinger operator $S_{\beta,k}$ on the half-line as the
scaling limit of $\beta$-ensembles at a soft edge where the equilibrium density vanishes to the order $2k+1/2$. We
show that the right tail of minus its lowest eigenvalue decays like $\exp(-c\,s^{(4k+3)/(4k+2)})$, with the exact
constant given by a weighted Gagliardo--Nirenberg inequality; the exponent was known from a lower bound of
Baslingker. For unitary ensembles the largest eigenvalue at such an edge has a much lighter tail, at most of order
$\exp(-c\,s^{2k+3/2})$, whenever the limiting edge kernel decays accordingly, so the proposed limit fails for every
scale constant. Baslingker proved this for $k=1$. For every $k\ge2$ it applies to an even polynomial potential
studied by Dai, Du and Lu, whose edge kernel is a dilation of the kernel of Claeys, Its and Krasovsky, so the
conjecture fails there.
\end{abstract}

\noindent\textbf{Keywords:} stochastic Airy operator; beta ensembles; higher-order Tracy--Widom laws; singular
edge; large deviations.

\noindent\textbf{MSC 2020:} 60B20; 15B52; 47B80.

\section{Introduction}\label{sec:intro}

Let $V$ be a real polynomial of even degree with positive leading coefficient, and consider $n$ points on $\R$ with
joint density proportional to
\begin{equation}\label{eq:ensemble}
\prod_{i<j}|x_i-x_j|^{\beta}\exp\Big(-\frac{\beta n}{2}\sum_{i=1}^n V(x_i)\Big).
\end{equation}
For $\beta=2$ this is the eigenvalue density of the unitary ensemble with weight $e^{-n\operatorname{Tr}V(M)}$. The
empirical measure converges to the equilibrium measure $\mu_V$ of $V$. At a right edge $E$ of its support the density
of $\mu_V$ generically vanishes like a square root, and the largest point, rescaled by $n^{2/3}$, converges to the
Tracy--Widom law; for every $\beta$ the limit is minus the lowest eigenvalue of the stochastic Airy operator
\cite{RRV,KRV}. For a real analytic $V$ the density of $\mu_V$ can also vanish at an edge like $(E-x)^{2k+1/2}$ for an
integer $k\ge1$. Krishnapur, Rider and Vir\'ag \cite[Conjecture~45]{KRVarxiv} (see also \cite[Conjecture~5]{Virag})
proposed that in this case the tridiagonal model of \eqref{eq:ensemble}, rescaled by $n^{2/(4k+3)}$, converges to
\begin{equation}\label{eq:operator}
S_{\beta,k}=-\frac{d^2}{dt^2}+t^{1/(2k+1)}+\frac{2}{\sqrt\beta}\,t^{-k/(2k+1)}\,W'(t),\qquad t>0,
\end{equation}
with a Dirichlet condition at $0$, where $W$ is a standard Brownian motion. In particular the lowest eigenvalue
$\Lambda_0$ of $S_{\beta,k}$ would be the limit in law of $\gamma n^{2/(4k+3)}(E-\lambda_{\max})$ for a constant
$\gamma>0$ determined by $V$.

Baslingker \cite{Bas} showed that this fails for $\beta=2$ and $k=1$. He proved the lower bound
$\Prob(-\Lambda_0>s)\ge\exp(-c_\beta s^{(4k+3)/(4k+2)})$ for all $k\ge1$ and compared it with the right tail
$\exp(-cs^{7/2})$ of the limiting law of the largest eigenvalue at $k=1$ \cite[(16)]{Bas}, which follows from the
edge kernel of Claeys and Vanlessen \cite{CV} and the tail estimate of Claeys, Its and Krasovsky \cite{CIK}. He noted
\cite[Remark~3]{Bas} that a heuristic tail estimate of Majumdar and Schehr \cite{MS} for $k>1$, together with his lower
bound, suggests that the conjecture does not hold in general.

Our first result identifies the right tail of $-\Lambda_0$ exactly. Put
\begin{equation}\label{eq:gd}
g=\frac{k}{2k+1},\qquad d=\frac1{2k+1},\qquad
v_k=\sup_{0\ne u\in C_c^\infty(0,\infty)}\frac{\int_0^\infty t^{-2g}u(t)^4\,dt}{\bigl(\int_0^\infty(u'(t)^2+u(t)^2)\,dt\bigr)^2}.
\end{equation}
\Cref{lem:vk} shows that $0<v_k<\infty$.

\begin{theorem}\label{thm:tail}
Let $\beta>0$ and $k\ge1$, and let $\Lambda_0$ be the lowest eigenvalue of $S_{\beta,k}$. Then
\[
\lim_{s\to\infty}s^{-(4k+3)/(4k+2)}\log\Prob(\Lambda_0<-s)=-\frac{\beta}{8v_k}.
\]
\end{theorem}

The exponent $(4k+3)/(4k+2)=\frac32-g$ is the one in Baslingker's lower bound \cite[(22)]{Bas}. Formally setting
$k=0$ in \eqref{eq:gd} gives the unweighted constant; on the half-line it is approached by translates of $u(t)=\operatorname{sech}t$, so that
$v_0=3/16$ and $\beta/(8v_0)=2\beta/3$, which is the right-tail constant of the Tracy--Widom law
$\Prob(-\Lambda_0>s)=\exp(-(2\beta/3+o(1))s^{3/2})$ as $s\to\infty$ for the stochastic Airy operator \cite{DV,RRV}. The proof of
\cref{thm:tail} does not cover $k=0$; see \cref{rem:k0}.

Our second result is the comparison at $\beta=2$, separated from its kernel input. Here and below $c$ and $C$
denote positive constants that may depend on $k$, $\beta$ and $V$ but not on $n$, $s$ or $u$; their values may
change from line to line. Let $K_n(x,y)=e^{-n(V(x)+V(y))/2}\sum_{j<n}\varphi_j(x)\varphi_j(y)$ be the
correlation kernel of \eqref{eq:ensemble} with $\beta=2$, where $\varphi_j$ are the orthonormal polynomials for the weight
$e^{-nV}$ on $\R$; $K_n(x,x)$ is the one-point density of the $n$ points.

\begin{theorem}\label{thm:counter}
Let $k\ge1$ and let $V$ be a real polynomial whose equilibrium measure has density $\psi_V$ with
$\psi_V(x)=c(E-x)^{2k+1/2}(1+O(E-x))$ as $x\uparrow E$, where $E$ is the right end of the support and $c>0$. Put
$p=2k+3/2$ and $h=1/p=2/(4k+3)$. Suppose that for some constant $\kappa>0$, depending only on $V$, the kernels
$\widetilde K_n(u,v)=\kappa^{-1}n^{-h}K_n(E+\kappa^{-1}n^{-h}u,E+\kappa^{-1}n^{-h}v)$ converge uniformly on compact subsets of
$\R^2$ to a continuous kernel $K$, and that
\begin{equation}\label{eq:H}
\int_s^\infty K(u,u)\,du\le C e^{-cs^{p}}\qquad(s\ge1)
\end{equation}
for some $c,C>0$. Then for $\beta=2$ there is no constant $\gamma>0$ for which $\gamma n^{h}(E-\lambda_{\max})$
converges in law to $\Lambda_0$, the lowest eigenvalue of $S_{2,k}$. The same holds along any subsequence of $n$.
\end{theorem}

\begin{corollary}\label{cor:all-k}
Let $k\ge2$ and let $V$ be the potential of \cite[Remark~1.7]{DDL}, recalled in \cref{sec:ensemble}, whose equilibrium
density vanishes like $c(a-x)^{2k+1/2}$ at the right end $E=a$ of its support. By \cite[Theorem~1.9]{DDL} and
\cite[(1.23)]{CIK}, the kernels of $V$ satisfy the hypotheses of \cref{thm:counter}. Hence the conjecture
\cite[Conjecture~45]{KRVarxiv} fails for $V$ at $\beta=2$, for every choice of the scale constant.
\end{corollary}

For $k\ge2$ the limit kernel at the edge of this potential is given by Dai, Du and Lu \cite[Theorem~1.9]{DDL}, a
preprint. It is a dilation of the kernel of Claeys, Its and Krasovsky \cite{CIK}, whose decay \cite[(1.23)]{CIK} gives
\eqref{eq:H} (\cref{sec:ensemble}). For general potentials with such an edge the limit kernel exists by
\cite[Proposition~1.7 and Appendix~B]{CKLW}; the edge model problem is also studied in \cite{Atkin}.
\Cref{thm:counter} isolates exactly what is used. For $k=1$ the limit kernel was identified
by Claeys and Vanlessen \cite{CV}, and it satisfies the pointwise bound $K(u,u)\le Ce^{-cu^{7/2}}$ for $u\ge1$ by
\cite[(1.23)]{CIK} and continuity, which implies \eqref{eq:H}; so \cref{thm:counter} gives Baslingker's theorem \cite[Theorem~4]{Bas}.

\Cref{thm:counter} uses only two features of $\Lambda_0$: that $-\Lambda_0$ has a right tail heavier than
$e^{-cs^{p}}$, and that it is a limit in law. Any proposed limit whose right tail is at least $e^{-Cs^2}$, for instance
the lowest eigenvalue of any semibounded operator whose quadratic form on one fixed test vector is a nondegenerate
Gaussian, is excluded in the same way. The theorem says nothing about what the correct limit is. The notation is collected in \cref{app:notation}.

\section{The operator and the proof of \texorpdfstring{\cref{thm:tail}}{Theorem 1.1}}\label{sec:operator}

Throughout, $k\ge1$ is fixed and $g,d$ are as in \eqref{eq:gd}, so $0<g<1/2$ and $d>0$. Write $\|\cdot\|$ for the norm
of $L^2(0,\infty)$, and for $u\in C_c^\infty(0,\infty)$ put
\[
\mathcal E(u)=\int_0^\infty(u'^2+u^2),\qquad P(u)=\int_0^\infty t^d u^2,\qquad X(u)=\int_0^\infty t^{-g}u(t)^2\,dW(t).
\]
For fixed $u$ the variable $X(u)$ is centred Gaussian with variance $\int t^{-2g}u^4$, and the covariances are
$\operatorname{Cov}(X(u_1),X(u_2))=\int t^{-2g}u_1^2u_2^2$.

\begin{lemma}\label{lem:vk}
$0<v_k<\infty$. More precisely $\int_0^\infty t^{-2g}u^4\le\bigl(1+\frac{1}{1-2g}\bigr)\mathcal E(u)^2$ for
$u\in C_c^\infty(0,\infty)$.
\end{lemma}

\begin{proof}
Positivity is clear. For $u\in C_c^\infty(0,\infty)$, $u(t)^2=-\int_t^\infty 2uu'$ gives
$\|u\|_\infty^2\le2\|u\|\,\|u'\|\le \mathcal E(u)$. Hence
$\int_0^1t^{-2g}u^4\le\|u\|_\infty^4/(1-2g)\le \mathcal E(u)^2/(1-2g)$, and
$\int_1^\infty t^{-2g}u^4\le\|u\|_\infty^2\|u\|^2\le \mathcal E(u)^2$.
\end{proof}

\subsection{The noise term}

For $j\ge0$ let $I_j=[j,j+1]$ and $Z_j(t)=\int_j^t r^{-g}\,dW(r)$ for $t\in I_j$; put
$B=\sup_{j\ge0}\sup_{t\in I_j}|Z_j(t)|$.

\begin{lemma}\label{lem:B}
$\E B<\infty$, and almost surely, for every $u\in C_c^\infty(0,\infty)$,
\begin{equation}\label{eq:Xbound}
|X(u)|\le B\bigl(\|u\|^2+4\|u\|\,\|u'\|\bigr)\le 3B\,\mathcal E(u).
\end{equation}
\end{lemma}

\begin{proof}
$Z_j$ is a continuous martingale on $I_j$ with quadratic variation at most $q_0=1/(1-2g)$ for $j=0$ and at most
$q_j=j^{-2g}$ for $j\ge1$. By the reflection principle for a time-changed Brownian motion,
$\Prob(\sup_{I_j}|Z_j|>b)\le4e^{-b^2/(2q_j)}$, so for $b\ge1$
\[
\Prob(B>b)\le4e^{-(1-2g)b^2/2}+4\sum_{j\ge1}e^{-b^2j^{2g}/2}\le4e^{-(1-2g)b^2/2}+4e^{-b^2/4}\sum_{j\ge1}e^{-j^{2g}/4},
\]
using $b^2j^{2g}/2\ge b^2/4+j^{2g}/4$ for $b,j\ge1$. The last sum is finite because $g>0$, so $\E B<\infty$.

Integrating by parts on each $I_j$ (Itô's product rule; $u^2$ is deterministic and smooth),
\begin{equation}\label{eq:ibp}
\int_{I_j}t^{-g}u^2\,dW=Z_j(j+1)\,u(j+1)^2-\int_{I_j}Z_j\,(u^2)'\,dt .
\end{equation}
We take the sum over $j$ of the right side of \eqref{eq:ibp} as the definition of $X(u)$ for all $u$ at once; for
each fixed $u$ it agrees almost surely with the It\^o integral, so the Gaussian description above is unchanged, and
$X$ is a quadratic form in $u$ on every path. By \eqref{eq:ibp},
$|\int_{I_j}t^{-g}u^2\,dW|\le B\,u(j+1)^2+2B\int_{I_j}|uu'|$. Writing $u(j+1)^2=u(x)^2+\int_x^{j+1}(u^2)'$ and
averaging over $x\in I_j$ gives $u(j+1)^2\le\int_{I_j}u^2+2\int_{I_j}|uu'|$. Summing over $j$ (finitely many terms are
nonzero) and using Cauchy--Schwarz gives the first inequality in \eqref{eq:Xbound}; the second follows from
$4\|u\|\|u'\|\le2(\|u\|^2+\|u'\|^2)$.
\end{proof}

Since \eqref{eq:Xbound} holds for all $u$ on one event of probability one, polarisation shows that $X$ extends
almost surely to a continuous quadratic form on the Sobolev space $H_0^1(0,\infty)$, and for every $\varepsilon>0$,
$|X(u)|\le\varepsilon\|u'\|^2+(B+4B^2/\varepsilon)\|u\|^2$. Let $\mathcal H$ be the completion of
$C_c^\infty(0,\infty)$ in the norm $\bigl(\|u'\|^2+\|u\|^2+P(u)\bigr)^{1/2}$. By the KLMN theorem
\cite{RS2}, the form
\[
q(u)=\|u'\|^2+P(u)+\tfrac{2}{\sqrt\beta}X(u),\qquad u\in\mathcal H,
\]
is almost surely closed and bounded below, with $C_c^\infty(0,\infty)$ as a form core, and it defines $S_{\beta,k}$.
Since $t^d\to\infty$, $\mathcal H$ embeds compactly into $L^2(0,\infty)$, so $S_{\beta,k}$ has discrete spectrum
bounded below, and
\begin{equation}\label{eq:var}
\Lambda_0=\inf_{0\ne u\in D}\frac{q(u)}{\|u\|^2}
\end{equation}
for any countable set $D\subset C_c^\infty(0,\infty)$ dense in $\mathcal H$, by the core property and the continuity of
$q$ on $\mathcal H$. This is the definition of \cite{KRVarxiv,Bas} through the quadratic form.

\subsection{Scaling}

For $s>0$ let $(U_sf)(t)=s^{1/4}f(\sqrt s\,t)$, a unitary map of $L^2(0,\infty)$ onto itself that maps
$C_c^\infty(0,\infty)$ onto itself and is continuous on $\mathcal H$. Then
\[
\|(U_sf)'\|^2=s\|f'\|^2,\qquad P(U_sf)=s^{-d/2}P(f),
\]
\[
\operatorname{Cov}\bigl(X(U_sf_1),X(U_sf_2)\bigr)=s^{1/2+g}\int_0^\infty t^{-2g}f_1^2f_2^2,
\]
the last by the change of variables $r=\sqrt s\,t$. Hence the centred Gaussian processes $(X(U_sf))_f$ and
$(s^{1/4+g/2}X(f))_f$, indexed by $C_c^\infty(0,\infty)$, have the same finite-dimensional distributions. Choose $D$ in
\eqref{eq:var} with $U_sD$ also dense in $\mathcal H$ (for example $D=U_s^{-1}D'$ with $D'$ countable and dense).
Then \eqref{eq:var} and $\|U_sf\|=\|f\|$ give
\[
\{\Lambda_0<-s\}=\Big\{\exists f\in D:\ s\|f'\|^2+s^{-d/2}P(f)+\tfrac{2}{\sqrt\beta}X(U_sf)<-s\|f\|^2\Big\},
\]
and since the event depends only on countably many coordinates of the process,
\[
\Prob(\Lambda_0<-s)=\Prob\Big(\exists f\in D:\ -X(f)>\tfrac{\sqrt\beta}{2}\bigl(s^{3/4-g/2}\mathcal E(f)
+s^{-d/2-1/4-g/2}P(f)\bigr)\Big).
\]
All three terms are homogeneous of degree two in $f$, so we may restrict to $\mathcal E(f)=1$. Write
$r_s=\frac{\sqrt\beta}{2}s^{3/4-g/2}$; then $r_s^2/(2v)=\beta s^{\frac32-g}/(8v)$ for every $v>0$, and
$\frac32-g=(4k+3)/(4k+2)$.

\begin{proof}[Proof of \cref{thm:tail}: upper bound]
Dropping $P(f)\ge0$,
\[
\Prob(\Lambda_0<-s)\le\Prob\Big(\sup_{f\in D,\,\mathcal E(f)=1}-X(f)>r_s\Big).
\]
The family $\{-X(f):f\in D,\,\mathcal E(f)=1\}$ is a countable centred Gaussian family. By \cref{lem:B} its supremum $Y$
satisfies $\E Y\le3\E B<\infty$, and by \cref{lem:vk} and the density of $D$ its largest variance is
$\sup_{\mathcal E(f)=1}\int t^{-2g}f^4=v_k$. The Borell--TIS inequality \cite{Borell,TIS}, \cite[Theorem~2.1.1]{AT} gives
\[
\Prob(Y>r)\le\exp\bigl(-(r-\E Y)^2/(2v_k)\bigr)\qquad(r>\E Y).
\]
With $r=r_s$ this gives
$\limsup_{s\to\infty}s^{-(3/2-g)}\log\Prob(\Lambda_0<-s)\le-\beta/(8v_k)$.
\end{proof}

\begin{proof}[Proof of \cref{thm:tail}: lower bound]
Fix $f\in C_c^\infty(0,\infty)$ with $\mathcal E(f)=1$ and $v=\int t^{-2g}f^4>0$. Taking this single $f$ in the scaled
event (we may include it in $D$),
\[
\Prob(\Lambda_0<-s)\ge\Prob\bigl(-X(f)>r_s(1+s^{-1-d/2}P(f))\bigr),
\]
and $-X(f)$ is centred Gaussian with variance $v$. Since $r_s(1+s^{-1-d/2}P(f))=r_s(1+o(1))$ as $s\to\infty$, the Gaussian tail gives
$\liminf_{s\to\infty}s^{-(3/2-g)}\log\Prob(\Lambda_0<-s)\ge-\beta/(8v)$. Letting $v$ increase to $v_k$ over such $f$
proves the lower bound.
\end{proof}

\begin{remark}\label{rem:k0}
Both $g<1/2$ and $g>0$ are used: the first for integrability of $t^{-2g}$ at the origin in \cref{lem:vk,lem:B}, the
second for the summability over the infinitely many intervals $I_j$ in \cref{lem:B}. At $g=0$ the supremum of
$-X(f)$ over $\mathcal E(f)=1$ is almost surely infinite once the confining term is dropped, so the argument above does not
apply to the stochastic Airy operator, although the constant it produces agrees with the known one.
\end{remark}

\section{Ensembles with an edge of exponent \texorpdfstring{$2k+1/2$}{2k+1/2}}\label{sec:ensemble}

For $k\ge1$ let $a>0$ be given by $a^{4k+2}=2^{2k+2}(2k+1)!/(4k+1)!!$ and let
\[
V(x)=\sum_{j=0}^{2k}(-1)^j\binom{2k+\frac12}{j}\frac{a^{2j}}{4k+2-2j}\,x^{4k+2-2j},
\]
the potential of \cite[Remark~1.7]{DDL}. It is even, of degree $4k+2$, with leading coefficient $1/(4k+2)$. Let
$Q(z)=(z^2-a^2)^{2k+1/2}$ on $\C\setminus[-a,a]$, with the branch for which $Q(z)/z^{4k+1}\to1$ as $z\to\infty$. Expanding
$Q(z)=z^{4k+1}(1-a^2z^{-2})^{2k+1/2}$ shows that $V'$ is the polynomial part of the Laurent expansion of $Q$ at infinity,
so $Q-V'=O(z^{-1})$. Dai, Du and Lu \cite[Remark~1.7]{DDL} state that $V$ satisfies \cite[Assumptions~1.1]{DDL} and
that its equilibrium density is the function $\rho$ below. We record the computation in the normalisation of
\eqref{eq:ensemble}, and in \cref{lem:ddl} we give deformation directions $V_j$ for which their condition (iii) holds;
the remark does not specify them.

\begin{lemma}[{\cite[Remark~1.7]{DDL}}]\label{lem:eq}
In the convention of \eqref{eq:ensemble} with $\beta=2$, the equilibrium measure of $V$ is $\rho(x)\,dx$ with
$\rho(x)=\frac1{2\pi}(a^2-x^2)^{2k+1/2}\mathbf 1_{[-a,a]}(x)$, where $\mathbf 1_{[-a,a]}$ is the indicator function of
$[-a,a]$. In particular $E=a$ and $\psi_V(x)=\frac{(2a)^{2k+1/2}}{2\pi}(a-x)^{2k+1/2}(1+O(a-x))$ as $x\uparrow a$.
\end{lemma}

\begin{proof}
Since $\Gamma(2k+3/2)=(4k+1)!!\sqrt\pi/2^{2k+1}$,
\[
\int\rho=\frac{a^{4k+2}}{2\pi}\int_{-1}^1(1-y^2)^{2k+1/2}\,dy=\frac{a^{4k+2}}{2\pi}\cdot\frac{\sqrt\pi\,\Gamma(2k+3/2)}{\Gamma(2k+2)}
=\frac{a^{4k+2}(4k+1)!!}{2^{2k+2}(2k+1)!}=1.
\]
Let $G(z)=\int\rho(x)(z-x)^{-1}dx$. The function $F=\frac12(V'-Q)$ is analytic on $\C\setminus[-a,a]$ and tends to $0$
at infinity. On $(-a,a)$ the boundary values of $Q$ from above and below are $Q_\pm(x)=\pm i(a^2-x^2)^{2k+1/2}$, so
$F_--F_+=\frac12(Q_+-Q_-)=2\pi i\rho(x)$, the same jump as $G$. So $F-G$ is analytic off $\{\pm a\}$, bounded near
$\pm a$, hence entire, and it tends to $0$ at infinity; thus $G=F$. In particular $F_++F_-=V'$ on $(-a,a)$, that is,
$V'(x)=2\,\mathrm{p.v.}\!\int\rho(t)(x-t)^{-1}dt$, and the effective potential
$\mathcal U(x)=V(x)-2\int\log|x-t|\rho(t)\,dt$ is constant on $[-a,a]$. For real $x\notin[-a,a]$, $\mathcal U'(x)=V'(x)-2G(x)=Q(x)$,
which is positive for $x>a$ and negative for $x<-a$. Hence $\mathcal U$ is strictly larger off $[-a,a]$ than on it, and the
Euler--Lagrange characterisation \cite[Theorem~I.3.3]{ST} identifies $\rho$ as the equilibrium density. The expansion
at $a$ follows from $a^2-x^2=(a-x)(2a-(a-x))$.
\end{proof}

So $V$ has one-interval support $[-a,a]$, positive density inside, and singular edges of exponent $2k+1/2$ at both
ends; only the right one enters \cref{thm:counter}. For $\beta=2$ the ensemble \eqref{eq:ensemble} with this $V$ is the
unitary ensemble $e^{-n\operatorname{Tr}V(M)}dM$ of \cite[(1.1)]{DDL} with all deformation parameters $t_j$ equal to $0$.

\begin{lemma}\label{lem:ddl}
$V$ satisfies \cite[Assumptions~1.1]{DDL}, with $A=a$ and $\psi\equiv1/(2\pi)$ in \cite[(1.40)]{DDL}, for suitable even
polynomials $V_1,\dots,V_{2k}$ with $\deg V_j=2j$.
\end{lemma}

\begin{proof}
Part (ii) of \cite[Assumptions~1.1]{DDL} is \cref{lem:eq}; the strict inequality off $[-a,a]$ that it requires
is shown in the proof of \cref{lem:eq}. In part (i), $V$ is an even polynomial, and since
$\deg V_j\le4k<\deg V$ and the leading coefficient of $V$ is positive, the growth condition \cite[(1.2)]{DDL} holds
uniformly for $|t_j|\le1$. It remains to find $V_j$ satisfying \cite[(1.41)]{DDL}.

Let $R(z)=((z-a)(z+a))^{1/2}$ on $\C\setminus[-a,a]$ with $R(z)/z\to1$ as $z\to\infty$. For an even polynomial $J$ the function
$RJ'$ is analytic off $[-a,a]$ and even, so at infinity $RJ'=\omega_J+O(z^{-2})$ with $\omega_J$ an even polynomial.
Moving the contour of \cite[(1.42)]{DDL} to a large circle shows that the function defined there, for $J=V_j$, is
$\omega_{V_j}$. Since $R(z)=z(1-a^2z^{-2})^{1/2}$, for $J=x^{2m}$ the polynomial $\omega_J$ has degree $2m$ and leading
coefficient $2m$. Hence $J\mapsto\omega_J$ is injective on the span of $x^2,\dots,x^{2j}$, and its image has dimension
$j$ inside the $(j+1)$-dimensional space of even polynomials of degree at most $2j$. For a polynomial $w$ let $\ell(w)$ be
the coefficient of $z^{-1}$ in the expansion of $w/R$ at infinity. Since $\omega_J/R=J'-(RJ'-\omega_J)/R=J'+O(z^{-3})$
and $J'$ is a polynomial, $\ell(\omega_J)=0$. As $\ell(1)=1$, the kernel of $\ell$ on even polynomials of degree at
most $2j$ has dimension $j$, so it is exactly the image of $J\mapsto\omega_J$.

For $m\ge0$, $\ell\bigl((z^2-a^2)^m\bigr)$ is the coefficient of $z^{-1}$ in
$(z^2-a^2)^{m-1/2}=z^{2m-1}(1-a^2z^{-2})^{m-1/2}$, that is, $\binom{m-1/2}{m}(-a^2)^m$, which is not $0$ because
$m-1/2$ is not an integer. Put
\[
\Pi_j(z)=(z^2-a^2)^{j-1}\Bigl(\ell\bigl((z^2-a^2)^j\bigr)-\ell\bigl((z^2-a^2)^{j-1}\bigr)(z^2-a^2)\Bigr),\qquad j=1,\dots,2k.
\]
Then $\Pi_j$ is even of degree $2j$ and $\ell(\Pi_j)=0$, so $\Pi_j=\omega_{V_j}$ for an even polynomial $V_j$ of
degree $2j$. The second factor of $\Pi_j$ takes the value $\ell((z^2-a^2)^j)\ne0$ at $\pm a$, so $\Pi_j$ vanishes at
$\pm a$ to order exactly $j-1$, which is \cite[(1.41)]{DDL}.
\end{proof}

The deformation directions $V_j$ enter only through $V+\sum_jt_jV_j$, and below all $t_j$ are $0$.

\begin{proof}[Proof of \cref{cor:all-k}]
Take $t_j=0$ for every $j$ and every $n$. Then all the limits $\widehat t_j$ of \cite[(1.63)]{DDL} are $0$, so all the
parameters of the limit kernel in \cite[Theorem~1.9]{DDL} are $0$, and $K_n$ is the kernel $K_n^{(\mathbf t)}$ of
\cite[(1.3)]{DDL} at $\mathbf t=0$. By \cref{lem:ddl} and \cite[Theorem~1.9]{DDL}, the relation \cite[(1.24)]{DDL}
holds at the right edge $a$, uniformly for $u$ and $v$ in compact subsets of $\R$. So, with $E=a$ and $\kappa$ the scaling
constant of \cite[Theorem~1.9]{DDL}, the kernels $\widetilde K_n$ of \cref{thm:counter} converge uniformly on compact
subsets of $\R^2$ to the kernel of \cite[(1.25)]{DDL},
\[
K(u,v)=\frac{\Psi_1(u)\Psi_2(v)-\Psi_1(v)\Psi_2(u)}{-2\pi i(u-v)},
\]
where $\Psi$ is the solution of the Riemann--Hilbert problem of \cite[Section~1.4]{DDL} with all
parameters equal to $0$, and $\Psi_1,\Psi_2$ are the analytic extensions of its entries $\Psi_{11},\Psi_{21}$ from the sector
between $\Gamma_1$ and $\Gamma_2$. The kernel $K$ is continuous.

This kernel is a dilation of the kernel of Claeys, Its and Krasovsky. Let $\nu=2^{2/(4k+3)}$, so that
$\nu^{(4k+3)/2}=2$, and put $\Phi(\zeta)=\nu^{-\sigma_3/4}\Psi(\zeta/\nu)$, where $\sigma_3=\mathrm{diag}(1,-1)$. The contour $\Gamma_1\cup\dots\cup\Gamma_4$ of
\cite[Section~1.4]{DDL}, with its orientation, and the jump matrices \cite[(1.33)--(1.35)]{DDL} are those of
\cite[(1.16)--(1.18)]{CIK}; the contour is invariant under $\zeta\mapsto\zeta/\nu$, and $\Phi$ is bounded near $0$. With
all parameters $0$ the exponent in \cite[(1.38)]{DDL} is $\frac{4}{4k+3}\zeta^{(4k+3)/2}$, so
$\frac4{4k+3}(\zeta/\nu)^{(4k+3)/2}=\frac2{4k+3}\zeta^{(4k+3)/2}$, the exponent of \cite[(1.21)]{CIK} with
$t_0=\dots=t_{2k-1}=0$. Since $(\zeta/\nu)^{-\sigma_3/4}=\nu^{\sigma_3/4}\zeta^{-\sigma_3/4}$, the expansion
\cite[(1.36)]{DDL} becomes
\[
\Phi(\zeta)=\zeta^{-\sigma_3/4}\mathcal N\bigl(I+\Phi_\infty\zeta^{-1/2}+O(\zeta^{-1})\bigr)e^{-\vartheta(\zeta)\sigma_3},\qquad
\vartheta(\zeta)=\frac{2}{4k+3}\zeta^{(4k+3)/2},
\]
with a constant matrix $\Phi_\infty$ and the matrix $\mathcal N$ of \cite[(1.37)]{DDL}, which is that of \cite[(1.20)]{CIK}. So
$\Phi$ solves the Riemann--Hilbert problem of \cite[Section~1]{CIK} with $t_0=\dots=t_{2k-1}=0$. By
\cite[Remark~1.4]{CIK} the kernel $K^{(k)}$ of \cite[(1.14)]{CIK} does not depend on which solution is used, so it may be
computed from $\Phi$. The sector between $\Gamma_1$ and $\Gamma_2$ is invariant under the dilation, and
$\Phi_{11}(\zeta)=\nu^{-1/4}\Psi_{11}(\zeta/\nu)$, $\Phi_{21}(\zeta)=\nu^{1/4}\Psi_{21}(\zeta/\nu)$. Comparing
\cite[(1.14)]{CIK} with the formula for $K$ above gives
\[
K(u,v)=\nu\,K^{(k)}(\nu u,\nu v;0,\dots,0).
\]

By \cite[(1.23)]{CIK} there are $c,C,x_0>0$ with $|K^{(k)}(x,y;0,\dots,0)|\le Ce^{-cx^{p}-cy^{p}}$ for $x,y\ge x_0$,
$x\ne y$; letting $y\to x$ and using continuity, the same bound holds for $x=y\ge x_0$. With $\nu^p=2$ this gives
$K(u,u)\le C\nu e^{-4cu^{p}}$ for $u\ge x_0/\nu$, and $K(u,u)$ is bounded on $[1,x_0/\nu]$ by continuity. So
$K(u,u)\le C'e^{-c'u^{p}}$ for $u\ge1$, and for $s\ge1$, since $u^p-s^p\ge ps^{p-1}(u-s)\ge u-s$ for $u\ge s$,
\[
\int_s^\infty K(u,u)\,du\le C'e^{-c's^{p}}\int_s^\infty e^{-c'(u-s)}\,du=\frac{C'}{c'}e^{-c's^{p}},
\]
which is \eqref{eq:H}. \Cref{thm:counter} now applies to $V$ with $E=a$, which proves \cref{cor:all-k}.
\end{proof}

\section{Proof of \texorpdfstring{\cref{thm:counter}}{Theorem 1.2}}

Let $L_n=\kappa n^{h}(\lambda_{\max}-E)$, and suppose that $L_{n_j}$ converges in law along some subsequence to a
random variable $L$ with values in $\R$. For $s<T$, the portmanteau theorem for the open interval $(s,T)$ and the
bound of a probability by an expected count give, with $N_j(s,T)$ the number of $i$ with
$s<\kappa n_j^h(\lambda_i-E)<T$,
\[
\Prob(s<L<T)\le\liminf_j\Prob(s<L_{n_j}<T)\le\limsup_j\E N_j(s,T)=\int_s^TK(u,u)\,du,
\]
where the last step is $\E\#\{i:\lambda_i\in I\}=\int_IK_n(x,x)\,dx$, the change of variables to $\widetilde K_n$, and uniform
convergence on the compact interval $[s,T]$. The function $K(u,u)$ is nonnegative, as a limit of one-point densities.
Letting $T\to\infty$ and using \eqref{eq:H},
\begin{equation}\label{eq:upper}
\Prob(L>s)\le\int_s^\infty K(u,u)\,du\le C'e^{-c s^{p}}\qquad(s\ge1).
\end{equation}
Only the limit kernel on compact intervals enters; the finite-$n$ kernels are never integrated over an unbounded
range.

Now suppose $\gamma n_j^h(E-\lambda_{\max})\to\Lambda_0$ in law for some $\gamma>0$. Then $L_{n_j}\to L$ with
$L=-(\kappa/\gamma)\Lambda_0$ in law, and by \cref{thm:tail} with $\beta=2$,
\[
\Prob(L>s)=\Prob\bigl(\Lambda_0<-(\gamma/\kappa)s\bigr)\ge\exp\bigl(-C''s^{(4k+3)/(4k+2)}\bigr)
\]
for large $s$. Since $(4k+3)/(4k+2)<2<p$, this contradicts \eqref{eq:upper}. \qed

\medskip
The lower bound used here needs much less than \cref{thm:tail}: by \eqref{eq:var} with a single fixed test function,
$-\Lambda_0$ dominates a nondegenerate Gaussian variable, so $\Prob(-\Lambda_0>s)\ge e^{-Cs^2}$, which already
contradicts \eqref{eq:upper}.

\section*{Acknowledgments and funding}

This work was done by multiple AI agents, Claude (Anthropic) and Codex (OpenAI), connected to Math
Brain~\cite{MathBrainRepo} and working together inside the Brain framework~\cite{BrainFramework}, a shared, persistent
memory in which the agents record claims, proofs, failed attempts and methods, retrieve one another's results,
coordinate who works on what, and audit each other's work. For more on the Brain and Math Brain, see the Math Brain
paper~\cite{MathBrainPaper}, the Math Brain repository~\cite{MathBrainRepo} and the Brain framework
repository~\cite{BrainFramework}.

The research presented in this paper was supported by the European Research Council (ERC) under the European
Union's Horizon Europe research and innovation programme (grant agreement No.\ 101041711). It was also supported by the Simons
Foundation (as part of the Collaboration on the Mathematical and Scientific Foundations of Deep Learning), by Heights Labs, by
Convex Nexus Capital, and by the Israel Science Foundation (grants 2258/19 and 4101/25).

\section*{Declaration of competing interest}

The author holds equity in Heights Labs and Convex Nexus Capital and has a management role at Convex Nexus Capital.

\section*{Declaration of generative AI and AI-assisted technologies}

This paper is part of a broader research project on proving open problems with AI systems.
The proofs were developed with AI systems (OpenAI Codex and Anthropic Claude agents working under the author's
direction), which also drafted the text. The conjecture is due to Krishnapur, Rider and Vir\'ag \cite{KRVarxiv}, the
case $k=1$ and the lower-bound exponent are due to Baslingker \cite{Bas}, and the edge kernel inputs are those of
\cite{DDL,CIK}. The author directed the work and is responsible for the content of the paper and for decisions
concerning its dissemination.

\section*{Data and code availability}

No data or code were used. All proofs are in the paper.

\begin{appendices}
\section{Notation}\label{app:notation}

\begin{tabular}{@{}>{\raggedright\arraybackslash}p{0.22\textwidth}>{\raggedright\arraybackslash}p{0.49\textwidth}>{\raggedright\arraybackslash}p{0.2\textwidth}@{}}
Symbol & Meaning & First appears\\[2pt]
$\beta$, $n$, $V$ & inverse temperature, number of points, potential & \eqref{eq:ensemble}\\
$\mu_V$, $\psi_V$ & equilibrium measure and its density & \cref{sec:intro}\\
$E$ & right end of the support of $\mu_V$ & \cref{sec:intro}\\
$\lambda_i$, $\lambda_{\max}$ & the points of \eqref{eq:ensemble} and the largest one & \cref{sec:intro}\\
$S_{\beta,k}$, $W$ & the operator \eqref{eq:operator}, standard Brownian motion & \eqref{eq:operator}\\
$\Lambda_0$ & lowest eigenvalue of $S_{\beta,k}$ & \cref{sec:intro}\\
$\gamma$, $\kappa$ & scale constants of the conjecture and of the kernel limit & \cref{sec:intro}\\
$g$, $d$, $v_k$ & $k/(2k+1)$, $1/(2k+1)$, the constant in \eqref{eq:gd} & \eqref{eq:gd}\\
$K_n$, $\widetilde K_n$, $K$ & correlation kernel, its rescaling, its limit & \cref{sec:intro}\\
$\varphi_j$ & orthonormal polynomials for the weight $e^{-nV}$ & \cref{sec:intro}\\
$p$, $h$ & $2k+3/2$ and $1/p$ & \cref{thm:counter}\\
$c$, $C$ & positive constants, changing from line to line & \cref{sec:intro}\\
$\Prob$, $\E$ & probability and expectation & \cref{sec:intro}\\
$\mathcal E(u)$, $P(u)$, $X(u)$ & energy, confining term, noise term of the form & \cref{sec:operator}\\
$I_j$, $Z_j$, $B$, $b$ & unit intervals, local noise integrals, their supremum, a level & \cref{sec:operator}\\
$q$, $\mathcal H$, $D$ & quadratic form of $S_{\beta,k}$, its form domain, a dense set & \cref{sec:operator}\\
$U_s$, $r_s$, $f$ & unitary dilation, the threshold $\frac{\sqrt\beta}{2}s^{3/4-g/2}$, a test function & \cref{sec:operator}\\
$a$, $\rho$, $Q$, $\mathcal U$ & right end, density, $(z^2-a^2)^{2k+1/2}$, effective potential & \cref{sec:ensemble}\\
$G$, $F$ & Cauchy transform of $\rho$ and $\frac12(V'-Q)$ & \cref{sec:ensemble}\\
$R$, $\omega_J$, $\ell(w)$ & $((z-a)(z+a))^{1/2}$, polynomial part of $RJ'$, coefficient of $z^{-1}$ in $w/R$ & \cref{sec:ensemble}\\
$V_j$, $t_j$, $\Pi_j$ & deformation directions and parameters of \cite{DDL}, the polynomial $\omega_{V_j}$ & \cref{sec:ensemble}\\
$\Psi$, $\Psi_1,\Psi_2$ & model solution of \cite[Section~1.4]{DDL}, its first column & \cref{sec:ensemble}\\
$\nu$, $\sigma_3$, $\Phi$, $K^{(k)}$ & $2^{2/(4k+3)}$, $\mathrm{diag}(1,-1)$, $\nu^{-\sigma_3/4}\Psi(\zeta/\nu)$, kernel of \cite[(1.14)]{CIK} & \cref{sec:ensemble}\\
$\mathcal N$, $\Phi_\infty$, $\vartheta$ & constant matrices and exponent of the model problem & \cref{sec:ensemble}\\
$x_0$ & a threshold in \cite[(1.23)]{CIK} & \cref{sec:ensemble}\\
$L_n$, $N_j(s,T)$ & rescaled largest point, count of rescaled points in $(s,T)$ & proof of \cref{thm:counter}
\end{tabular}
\end{appendices}

\begin{thebibliography}{99}

\bibitem{AT}
R.~J.~Adler and J.~E.~Taylor, \emph{Random fields and geometry}, Springer Monographs in Mathematics, Springer, New
York, 2007, \doi{10.1007/978-0-387-48116-6}.

\bibitem{Atkin}
M.~R.~Atkin, \emph{A Riemann--Hilbert problem for equations of Painlev\'e type in the one matrix model with
semi-classical potential}, arXiv:1504.04539v1, 2015,
\url{https://arxiv.org/abs/1504.04539v1}.

\bibitem{Bas}
J.~Baslingker, \emph{Stochastic domination in beta ensembles}, J. Math. Phys.\ 66 (2025), no.~5, 053303,
\doi{10.1063/5.0254593}; numbering as in arXiv:2304.09236v2.

\bibitem{Borell}
C.~Borell, \emph{The Brunn--Minkowski inequality in Gauss space}, Invent. Math.\ 30 (1975), no.~2, 207--216,
\doi{10.1007/BF01425510}.

\bibitem{CIK}
T.~Claeys, A.~Its and I.~Krasovsky, \emph{Higher-order analogues of the Tracy--Widom distribution and the Painlev\'e
II hierarchy}, Comm. Pure Appl. Math.\ 63 (2010), no.~3, 362--412, \doi{10.1002/cpa.20284}; numbering as in arXiv:0901.2473v1.

\bibitem{CKLW}
T.~Claeys, A.~B.~J.~Kuijlaars, K.~Liechty and D.~Wang, \emph{Propagation of singular behavior for Gaussian
perturbations of random matrices}, Comm. Math. Phys.\ 362 (2018), 1--54, \doi{10.1007/s00220-018-3195-8}; numbering as in arXiv:1608.05870v1.

\bibitem{CV}
T.~Claeys and M.~Vanlessen, \emph{Universality of a double scaling limit near singular edge points in random matrix
models}, Comm. Math. Phys.\ 273 (2007), no.~2, 499--532, \doi{10.1007/s00220-007-0256-9}.

\bibitem{DDL}
D.~Dai, J.-H.~Du and C.~Lu, \emph{Asymptotics of Hankel determinants for potentials with singular edge points},
preprint, arXiv:2609.26412v1, 2026,
\url{https://arxiv.org/abs/2609.26412v1}.

\bibitem{DV}
L.~Dumaz and B.~Vir\'ag, \emph{The right tail exponent of the Tracy--Widom $\beta$ distribution}, Ann. Inst. Henri
Poincar\'e Probab. Stat.\ 49 (2013), no.~4, 915--933, \doi{10.1214/11-AIHP475}.

\bibitem{KRV}
M.~Krishnapur, B.~Rider and B.~Vir\'ag, \emph{Universality of the stochastic Airy operator}, Comm. Pure Appl. Math.\
69 (2016), no.~1, 145--199, \doi{10.1002/cpa.21573}.

\bibitem{MS}
S.~N. Majumdar and G.~Schehr, \emph{Top eigenvalue of a random matrix: large deviations and third order phase
transition}, J. Stat. Mech. Theory Exp. (2014), P01012, \doi{10.1088/1742-5468/2014/01/P01012}.

\bibitem{KRVarxiv}
M.~Krishnapur, B.~Rider and B.~Vir\'ag, \emph{Universality of the stochastic Airy operator}, arXiv:1306.4832v1, 2013,
\url{https://arxiv.org/abs/1306.4832v1}.

\bibitem{RRV}
J.~A.~Ram\'irez, B.~Rider and B.~Vir\'ag, \emph{Beta ensembles, stochastic Airy spectrum, and a diffusion}, J. Amer.
Math. Soc.\ 24 (2011), no.~4, 919--944, \doi{10.1090/S0894-0347-2011-00703-0}.

\bibitem{RS2}
M.~Reed and B.~Simon, \emph{Methods of modern mathematical physics II: Fourier analysis, self-adjointness}, Academic
Press, New York, 1975.

\bibitem{ST}
E.~B.~Saff and V.~Totik, \emph{Logarithmic potentials with external fields}, Grundlehren der mathematischen
Wissenschaften 316, Springer, Berlin, 1997.

\bibitem{TIS}
B.~S.~Tsirelson, I.~A.~Ibragimov and V.~N.~Sudakov, \emph{Norms of Gaussian sample functions}, in Proceedings of the
Third Japan--USSR Symposium on Probability Theory, Lecture Notes in Math.\ 550, Springer, Berlin, 1976, 20--41.

\bibitem{Virag}
B.~Vir\'ag, \emph{Operator limits of random matrices}, in Proceedings of the International Congress of Mathematicians,
Seoul 2014, Vol.~IV; arXiv:1804.06953.

\bibitem{MathBrainRepo}
Y.~Shmalo, \emph{Math Brain: public snapshot, code and tests}, 2026.
\url{https://github.com/yspennstate/math-brain-public}

\bibitem{MathBrainPaper}
Y.~Shmalo, \emph{Math Brain: shared, auditable memory for multi-agent mathematical research}, October 2026.
\url{https://github.com/yspennstate/math-brain-public/tree/main/papers/math-brain-2026-10-24},
\doi{10.5281/zenodo.24150837}.

\bibitem{BrainFramework}
Y.~Shmalo, \emph{Brain Framework}, public software repository, 2026.
\url{https://github.com/yspennstate/brain-framework-public}

\end{thebibliography}

\end{document}
